%% kapa: exam=1 round=24 year=2013 period=2 form=A subject=appraisal_law subject_ko=부동산관계법규 seq=13 number=53
%% source: Q-Net 감정평가사 기출문제 (cst00309) artlSeq=5073802 file=extracted/1008/5073802/2013년 제24회 감정평가사 1차시험 기출문제/2교시(A형).hwp
%% conversion: kordoc 4.24.51 -> tools/gen1.py (auto); review=proofread-sonnet; flags=none
\begin{question}{53}
국토의 계획 및 이용에 관한 법령상 용도지역, 용도지구, 용도구역의 지정에 관한 설명으로 옳지 않은 것은?
\begin{choices}
  \item 동일한 토지에 2개 이상의 용도지역을 중복하여 지정할 수 있다.
  \item 동일한 토지에 2개 이상의 용도지구를 중복하여 지정할 수 있다.
  \item 동일한 토지에 용도지역과 용도지구를 중복하여 지정할 수 있다.
  \item 동일한 토지에 용도구역과 용도지역을 중복하여 지정할 수 있다.
  \item 동일한 토지에 용도구역과 용도지구를 중복하여 지정할 수 있다.
\end{choices}
\end{question}
\answer{1}
